% ============================================================
% 欧拉《刚体运动理论》现代汉语选编与重述 —— 合册 PDF 导言区
% 体例：B5 tufte-book，正文窄栏 + 右边栏批注（楷体，不编号）
% 编译：xelatex（latexmk 驱动，见 build/book.sh）
% ============================================================
\documentclass[b5paper,oneside]{tufte-book}

% ---------- 中文（xeCJK；heading=false 避免与 tufte 标题冲突） ----------
\usepackage[UTF8,fontset=none,heading=false]{ctex}
\setCJKmainfont{Noto Serif SC}[AutoFakeBold=2.6,AutoFakeSlant=0.15]
\setCJKsansfont{SimHei}[AutoFakeBold=2.6]
\setCJKfamilyfont{zhkai}{KaiTi}[AutoFakeBold=2.2]
\setCJKfamilyfont{zhh}{SimHei}[AutoFakeBold=2.6]
\newcommand{\kaishu}{\CJKfamily{zhkai}}
\newcommand{\heiti}{\CJKfamily{zhh}}
% 公式中的中文（\mathrm{外}、公式内句号等）用文本字体排
\xeCJKsetup{CJKmath=true}

% ---------- 西文 ----------
\setmainfont{TeX Gyre Pagella}
\setsansfont{TeX Gyre Adventor}
\setmonofont{Consolas}[Scale=0.86]

% ---------- 数学与表格 ----------
\usepackage{amsmath,amssymb}
\usepackage{booktabs,longtable,array}
\usepackage{etoolbox}
% 图形搜索路径（xelatex 在仓库根目录运行，图形在 build/ 下）
\graphicspath{{build/}}
\AtBeginEnvironment{longtable}{\small}
% pandoc 片段里的紧凑列表
\providecommand{\tightlist}{%
  \setlength{\itemsep}{1pt}\setlength{\parskip}{0pt}}

% ---------- 标题体例（节与子节不自动编号，沿用底稿手写序号） ----------
\usepackage{titlesec}
\titleformat{\section}
  {\normalfont\zihao{4}\heiti}
  {}{0em}{}
\titleformat{\subsection}
  {\normalfont\zihao{-4}\heiti}
  {}{0em}{}
\titlespacing*{\section}{0pt}{1.5em}{0.65em}
\titlespacing*{\subsection}{0pt}{1.15em}{0.5em}

% ---------- 超链接（最后加载；bookmark 修正大纲顺序） ----------
\usepackage{xcolor}
\definecolor{eulerlink}{RGB}{18,64,132}
\definecolor{eulerurl}{RGB}{0,110,110}
\usepackage{hyperref}% 若已被更早载入则无操作
\hypersetup{colorlinks=true,linkcolor=eulerlink,urlcolor=eulerurl,
            citecolor=eulerlink,pdfstartview=FitH}
\usepackage{bookmark}
% 目录含节；书签全部由 TOC 自动生成，不用手工 \pdfbookmark（避免重复）
\setcounter{tocdepth}{1}

% ---------- 批注：复用 tufte sidenote 的边栏自动避让，隐藏编号 ----------
\setsidenotefont{\footnotesize\kaishu}
\renewcommand{\thefootnote}{}% 边栏批注不编号（计数器照常推进以保证避让）
\newcommand{\pnote}[2][批注]{\sidenote[]{\heiti【#1】\;\allowbreak #2}}

% ---------- 章与部 ----------
% \eulerchapter{部}{号}{题}{label} —— 分部内沿用底本章号
\newcommand{\eulerchapter}[4]{%
  \clearpage
  \setcounter{chapter}{#2}\addtocounter{chapter}{-1}%
  \refstepcounter{chapter}%
  \phantomsection\label{#4}%
  \vspace*{1.8em}%
  {\centering\zihao{-3}\heiti 第#2章\par}%
  \vspace{0.45em}%
  {\centering\zihao{3}\heiti #3\par}%
  \vspace{1.4em}%
  \markboth{第#2章\quad #3}{}%
  \addcontentsline{toc}{chapter}{第#2章　#3}%
}
% 卷首无编号章：\eulerfrontchapter{题}{label}
\newcommand{\eulerfrontchapter}[2]{%
  \clearpage
  \phantomsection\label{#2}%
  \vspace*{1.8em}%
  {\centering\zihao{3}\heiti #1\par}%
  \vspace{1.4em}%
  \markboth{#1}{}%
  \addcontentsline{toc}{chapter}{#1}%
}
% 章铭（英文原章题）
\newcommand{\eulerepigraph}[1]{%
  \begin{center}\itshape\small\color{black!62}#1\end{center}%
  \vspace{0.3em}}
% 节（入目录；书签由 hyperref 从 TOC 自动生成）
\newcommand{\esection}[2]{%
  \section*{#1}%
  \phantomsection\addcontentsline{toc}{section}{#1}}
\newcommand{\esubsection}[1]{\subsection*{#1}}

% ---------- 选编说明（章前框） ----------
\newenvironment{xuanbian}
  {\begin{quote}\small\kaishu{\heiti【选编说明】}\;}
  {\end{quote}}

% ---------- 插图：图题图下居中（1973 原版样式），全书连续编号 ----------
\newcommand{\figcap}[1]{%
  \refstepcounter{figure}%
  \par\nopagebreak\vspace{0.45em}%
  {\centering\footnotesize 图~\thefigure\quad #1\par}}
% 宽图：跨正文栏与边栏，随文就位（非浮动）
\newcommand{\figwide}[2]{%
  \par\medskip\noindent
  \begin{fullwidth}\centering
    \includegraphics[width=\linewidth]{#1}%
    \figcap{#2}%
  \end{fullwidth}
  \medskip\par\noindent}
% 窄图：正文栏内浮动居中
\newcommand{\fignorm}[2]{%
  \begin{figure}[!h]\centering
    \includegraphics[width=0.92\linewidth]{#1}%
    \figcap{#2}%
  \end{figure}}

% ---------- 部页：\eulerpart{第X部分}{名}{导语}{anchor} ----------
\newcommand{\eulerpart}[4]{%
  \clearpage
  \thispagestyle{empty}%
  \vspace*{0.30\textheight}%
  \begin{fullwidth}\centering
    {\zihao{1}\heiti #2\par}\vspace{1.2em}
    {\large\color{black!55}#1\par}\vspace{2.2em}
    {\small\color{black!55}#3\par}
  \end{fullwidth}%
  \clearpage}

% ---------- 版面微调 ----------
\linespread{1.06}
